% !TeX root = ../教材题库_全部条目_转写工作稿.tex
% Source: XJ-B1-C06-EX-P219-Y1028
% 第6章 / 6.2 抽样 / 教材 p.219 / PDF 第226页
% OCR未校对：公式与文字均可能有误；未完成数学TeX转写；禁止直接发给学生
% 原题及相邻原解截图（核对用）：../题目原图/XJ-B1-C06-EX-P219-Y1028-01.webp
% 本文件只在转写完成后作为题干使用；OCR 条目尚不是可靠题干。
\begin{bankitem}{6.2 例题 例题2}{第6章 / 6.2 抽样 / 教材 p.219 / PDF 第226页}{XJ-B1-C06-EX-P219-Y1028}{OCR未校对：公式与文字均可能有误；未完成数学TeX转写；禁止直接发给学生}
\begin{ocrdraft}
[教材p.219 / PDF第226页]
\par
00下列问题中，采用哪种抽样方法较为合理？
\par
（1）一个微波炉厂的质量检查组想了解某批次1000台微波炉的使用寿命．
\par
（2）每年6月6日是“全国爱眼日”．某县卫生部门要调查该县中小学生视力保
\par
护情况，已知该县现有在校小学生12000名、初中生10000名、高中生6000名．
\par
（3）立德中学要调查本校九年级4名学生的身高和体重情况，以供营养师参
\par
考进而指导食堂伙食营养搭配．
\par
解（1）由于总体容量较大，可采用随机数法进行抽样．
\par
（2）由于总体容量大，并且具有明显的层次性，因而应当先采用分层抽样，然
\par
后在每层中可采用随机数法进行抽样．
\par
（3）由于总体容量较大，且男女学生在身高和体重方面有较大的差异，所以应
\par
当先采用分层抽样，然后在每层中对男生和女生分别用抽签法进行抽样．
\par
科学的抽样方法可以提高用样本估计总体的精度，历史上曾发生由于采取不正确
\par
的抽样方法而导致调查结论严重失真的教训．但影响精度的不仅有抽样方法，还有样
\par
本容量．一般地，样本容量较大时，可以获得较精确的估计．
\end{ocrdraft}
\end{bankitem}
